% Introduction: research approach and contributions paragraph.
% Result sentence filled in from outputs/hem_vs_dmon111_n6 (changes-from-claude.md #12).
To answer this question, we compare three pooling regimes on the same task: a fixed HEM coarsening computed once from the PPI graph, a fully learned regime in which every pooling step is learned, and a hybrid regime in which the first $n$ levels are fixed and a single learned layer is applied to the resulting smaller graph. DiffPool and DMoN serve as the learned objectives. Applying DMoN to a graph of this size required several adaptations: a progressive cluster schedule and sparsified pooled adjacency matrices for the fully learned regime, and dropout on the cluster assignments for both. A central difficulty is that pooling choices also change how much information reaches the classifier: a shallow fixed coarsening passes close to one value per gene, whereas a learned layer compresses the graph into a few dozen clusters. We therefore compare fixed and learned clustering at matched depth, node count and classifier input size, so that the only difference is whether the final grouping is fixed or learned. Under this matched comparison, fixed HEM clustering outperformed learned DMoN clustering in every repetition ($93.2\%$ vs.\ $83.8\%$ mean test accuracy; $94.4\%$ vs.\ $89.4\%$ for the three-model ensembles) and was far more stable across repetitions, so the advantage of the fixed hierarchy is not explained by its wider bottleneck. Across our experiments, learning was most effective when confined to the small graph at the end of a fixed coarsening; learning every pooling step was both far more expensive and less accurate within our compute budget. Along the way, we found that comparisons of this kind are fragile: evaluation-scope and checkpointing errors changed reported accuracies, and results only became comparable after we adopted a shared test split, a deterministic evaluation order, softmax-averaged ensembling and matched classifier input sizes. The main contributions of this thesis are therefore (i) a controlled comparison of fixed, fully learned and hybrid pooling with two learned objectives on a genomic classification task, including a comparison of fixed and learned clustering at equal depth and bottleneck width, (ii) an adaptation of DMoN to a graph with over ten thousand nodes, and (iii) an evaluation protocol, together with a record of the pitfalls that motivated it.
